% ---
% id: foundations/req:rounding-direction
% kind: req
% title: "Rounding direction of emitted bounds"
% summary: >-
%   An emitted scalar bound is rounded onto the lattice in the direction that
%   weakens the claim it carries, never to nearest, and the direction is
%   recorded. A complementary quantity is recomputed from rounded weights and
%   never rounded on its own.
% prominence: standard
% status: accepted
% uses: [def:binary64, def:declared-data]
% mentions: []
% assumes: []
% refs: []
% provenance: newly-derived
% lean:
% ---

A lower bound rounded down, or an upper bound rounded up, remains a bound of the
same kind, so a directed rounding keeps what the bound certifies without a
rounding-error budget.

\begin{requirement}\label{req:rounding-direction}
(D-DIR) Every emitted scalar bound is rounded onto the lattice in the direction
that weakens the claim it carries, never to nearest: a lower bound \emph{down},
an upper bound \emph{up}, a constraint constant in the direction that enlarges
the feasible set for its declared relation. The direction is recorded per
bound, and where a family's completeness depends on a rounded weight the
complementary quantity is \emph{recomputed} from the rounded weights, never
rounded independently.

It binds every producer that declares a lattice exponent $k$
(Definition~\ref{def:declared-data}), for each scalar bound it emits, the
lattice being that of the lattice values at $k$
(Definition~\ref{def:binary64}\,(iv)); the claim that states the emission names
the direction of each such bound. A producer that declares no lattice exponent,
and that constructs an emitted scalar bound exactly and then rounds it to a
binary64 number, rounds it in the same direction, never to nearest.
\end{requirement}
